% !PRINTED VI
\weberpage{VI}
\noindent Wenn auch die Algebra zum Theil über die Zahlentheorie hinausgeht, und in andere Gebiete, z.~B. die Functionentheorie oder in ihren Anwendungen auch in die Geometrie hinüber greift, so ist doch die Zahlenlehre immer das vorzüglichste Beispiel für alle algebraischen Betrachtungen, und die Fragen der Zahlentheorie, die heute im Vordergrund des Interesses stehen, sind vorwiegend algebraischer Natur. Hierdurch war der Weg bezeichnet, den ich in meiner Arbeit zu gehen hatte.

Der grosse Stoff ist in zwei Bände vertheilt. Der erste Band enthält den elementaren Theil der Algebra, den man mit einem hergebrachten Ausdruck als \gesperrt{Buchstabenrechnung} bezeichnen kann, sodann die Vorschriften über die numerische Berechnung der Gleichungswurzeln und die Anfänge der \gesperrt{Galois}'schen Theorie.

Der zweite Band, der dem ersten hoffentlich in kurzer Zeit folgen wird, soll die allgemeine Theorie der endlichen Gruppen, die Theorie der linearen Substitutionsgruppen und Anwendungen auf verschiedene einzelne Probleme bringen, und soll abschliessen mit der Theorie der algebraischen Zahlen, wo der Versuch gemacht ist, die verschiedenen Gesichtspunkte, unter denen diese Theorie bisher betrachtet worden ist, zu vereinigen.

Endlich soll der zweite Band ein alphabetisches Register über beide Bände bringen.

Wie es bei einer Disciplin, die in rascher Entwickelung begriffen ist, und an der von den verschiedensten Seiten gearbeitet wird, nicht anders sein kann, ist auch in der Algebra der Sprachgebrauch und die Bezeichnungsweise sehr mannigfaltig und häufig nicht übereinstimmend. Dadurch wird eine einheitliche Darstellung und das Eindringen in die verschiedenen Arbeiten sehr erschwert.

Ich habe mich daher bemüht, eine möglichst zweckmässige Ausdrucksweise einheitlich beizubehalten, und habe mich dabei vielfach mit Fachgenossen berathen. Ich darf die Hoffnung aussprechen, dadurch zur Befestigung einer einheitlichen Terminologie beigetragen zu haben.

Die Literaturnachweisungen und historischen Notizen, die in dem Buche gegeben sind, machen in keiner Weise den Anspruch auf Vollständigkeit, wenn ich auch bemüht gewesen bin,
